\section{Broader Impact}
This project aims to improve the empirical study of reward misspecification and
tampering by providing a compact benchmark with explicit assumptions and strict
artifact rules. If used well, such a benchmark can make safety discussions more
concrete: researchers can compare tasks, metrics, and monitor channels without
claiming more realism than the environments justify.

There are also risks. Compact safety benchmarks can be overinterpreted, used as
checklists for deployment readiness, or optimized against in a way that
encourages benchmark gaming rather than better safety science. We therefore
recommend treating \framework as a diagnostic tool for failure analysis and
method stress-testing, not as a certification suite.

Because the environments make corruption channels explicit, they may also serve
as toy examples of exploit strategies. We view that risk as limited by the
symbolic and low-stakes nature of the tasks, but it reinforces the importance of
presenting the benchmark as measurement infrastructure rather than a leaderboard.
